\documentclass[11pt]{article}
\usepackage[a4paper,margin=25mm]{geometry}
\usepackage{amsmath,amssymb,amsthm}
\usepackage[hidelinks]{hyperref}
\emergencystretch=2em
\newtheorem{theorem}{Theorem}
\title{A square-tessellation quotient below the claimed spectral bound\\
\large Disproof of Conjecture 00000007717}
\author{gaochengzhecpu}\date{}
\begin{document}\maketitle
\begin{abstract}
The square tessellation, quotiented by the period subgroup
$64\mathbb Z\times4\mathbb Z$, has vertex degree four and
four edges per cell. Its actual combinatorial Laplacian has
a positive eigenvalue $2-2\cos(\pi/32)<3/256$.
Consequently its first nonzero eigenvalue is below
$(3/16)/4^2$, contradicting the asserted constant $c=3/16$.
Lean verifies the quotient, connectivity, degree, cells,
real eigenvector and least positive eigenvalue.
\end{abstract}

\section*{The admissible tessellation and its quotient}
Take the ordinary planar square tessellation with vertex
set $\mathbb Z^2$ and edges joining points whose difference
is $(\pm1,0)$ or $(0,\pm1)$. The subgroup
$\Lambda=64\mathbb Z\times4\mathbb Z$ consists of periods
of this tessellation. Its quotient graph has vertex set
\[
 V=(\mathbb Z/64\mathbb Z)\times(\mathbb Z/4\mathbb Z).
\]
Put $a=(1,0)$ and $b=(0,1)$. A vertex $v$ is adjacent
exactly to $v+a,v-a,v+b,v-b$. These are four distinct
vertices, none equal to $v$, so the graph is simple and
has degree exactly four. Horizontal and vertical steps
connect every vertex to the origin. Every square cell
projects to the four distinct vertices
\[
 v,\quad v+a,\quad v+a+b,\quad v+b,
\]
with consecutive vertices adjacent, including the last to
the first. Thus the source's bounds hold with $D=4$ and
$E=4$.

The quotient map $\mathbb Z^2\to V$ reduces the two
coordinates modulo $64$ and $4$. It is surjective and its
kernel is exactly $\Lambda$. Two quotient vertices are
adjacent if and only if they have adjacent lifts. These
facts are formalized, so the construction is an actual
quotient of a tessellation 1-skeleton.
The statement does not require using the full group of
all translation symmetries or a smallest period cell.
Our chosen subgroup is a period subgroup, and no such
extra restriction is imposed here.

\newpage
\section*{An exact small eigenvalue}
For the combinatorial Laplacian, with no normalization,
\[
 (Lf)(i,j)=4f(i,j)-f(i+1,j)-f(i-1,j)
                   -f(i,j+1)-f(i,j-1).
\]
Let $\theta=\pi/32$ and $f(i,j)=\cos(i\theta)$.
This is well-defined modulo $64$, independent of $j$,
and nonzero since $f(0,0)=1$. The cosine addition formula
gives
\[
 Lf=(2-2\cos\theta)f.
\]
Write $\lambda=2-2\cos\theta$. Since $0<\theta<\pi$,
$\cos\theta<1$ and $\lambda>0$. The elementary inequality
$\cos x\geq1-x^2/2$ gives
\[
 \lambda\leq\theta^2=\frac{\pi^2}{1024}
        <\frac{12}{1024}=\frac3{256},
\]
where $\pi^2<12$ is certified exactly in Lean from a
rational upper bound for $\pi$.

The real Laplacian is positive semidefinite, so its
nonzero eigenvalues are positive. In this finite graph,
the set of positive eigenvalues is finite and nonempty,
as it contains $\lambda$. It has a least element
$\lambda_1$, and
\[
 0<\lambda_1\leq\lambda<\frac3{256}
       =\frac{3/16}{D^2}\qquad(D=4).
\]
\begin{theorem}
The constant $c=3/16$ asserted in Conjecture 00000007717
does not satisfy its claimed bound on all periodic
tessellation cell graphs.
\end{theorem}
\begin{proof}
The quotient above is connected, has $D=E=4$, and has
first nonzero eigenvalue strictly smaller than
$(3/16)/4^2$. It therefore violates that bound.
\end{proof}

This finite example refutes the explicit value claimed
for the optimal universal constant, hence the stated
conjunction. It does not rely on a disconnected graph,
weighted edges, growing degrees, or a normalized
Laplacian. No claim to formalize an entire family of
arbitrarily large quotients is needed.

\newpage
\section*{Formalization and reproducibility}
The graph is a Mathlib \texttt{SimpleGraph} on the actual
finite type \texttt{ZMod 64} times \texttt{ZMod 4}.
Theorems \texttt{project\_kernel} and
\texttt{quotient\_adjacency} identify the quotient and its
adjacency relation. Degree, connectivity and four distinct
cell vertices are checked separately.

The proof first uses the actual additive character on
$\mathbb Z/64\mathbb Z$ as a complex eigenvector. It
identifies the eigenvalue with the real number $\lambda$
and takes real parts to obtain a nonzero eigenvector
of the actual real Mathlib Laplacian matrix.
The theorem \texttt{spectral\_gap\_exists\_and\_small}
uses the matrix linear map's actual eigenvalues and
constructs their least positive element. The theorem
\texttt{counterexample} bundles the graph properties and
this strict bound. Positive semidefiniteness is also
verified. A small scalar estimate is therefore never
substituted for the spectral gap itself.

An independent exact-arithmetic check in
\texttt{verify.py} builds all 256 vertices and their
undirected edges. For the mean-zero triangular-wave
function $h(i,j)=\min(i,64-i)-16$, it obtains
\[
 \sum_v h(v)=0,\qquad \sum_v h(v)^2=21888,
 \qquad \sum_{\{u,v\}\in E}(h(u)-h(v))^2=256.
\]
Its Rayleigh quotient is $2/171<3/256$.
This is supplementary verification; Lean uses the
explicit trigonometric eigenvector instead.

Lean 4.19.0 and Mathlib revision\\
\texttt{c44e0c8ee63ca166450922a373c7409c5d26b00b}
are pinned. From the submitted \texttt{lean/} directory:
\begin{verbatim}
lake exe cache get
lake build
lake env lean -DwarningAsError=true Main.lean
\end{verbatim}
Run \texttt{python verify.py} and \texttt{tectonic main.tex}
for the auxiliary check and document. Fresh full build,
strict Lean, axiom audit, source provenance and PDF
verification records are included. Only standard
foundational axioms occur; no proof placeholders,
custom axioms or external computation oracles are used.

An agent developed the proof and supplied a self-review.
The parent separately reviewed its mathematical scope,
complete Lean source and final artifacts. No external
independent review or novelty is claimed for the elementary
Fourier calculation.

\begin{thebibliography}{9}
\bibitem{source} The Last Math Competition,
Conjecture 00000007717. The exact bilingual statement and
commit-pinned provenance are included as \texttt{SOURCE.md}
and its verification record.
\end{thebibliography}
\end{document}
